\chapter{Hepatitis E}
\index{Hepatitis E}
\label{sec:Hepatitis E}

\section{Overview}
Hepatitis E is an inflammatory liver disease caused by hepatitis E virus (HEV, Orthohepevirus A), a non-enveloped single-stranded RNA virus in the Hepeviridae family, with four major genotypes: genotypes 1 and 2 cause waterborne epidemics in developing countries, and genotypes 3 and 4 cause zoonotic infection (primarily from swine) in developed countries. An estimated 20 million HEV infections occur annually, with 3.3 million symptomatic cases and 44,000 deaths. This infection is transmitted through various routes depending on the specific pathogen. It remains a significant public health concern, particularly in regions with limited healthcare access. Infectious diseases are caused by pathogenic microorganisms including bacteria, viruses, fungi, and parasites that invade the body and multiply, leading to tissue damage and immune response. Transmission occurs through various routes: respiratory droplets, direct contact, contaminated food or water, vectors, or blood and body fluids. The incubation period varies from hours to years depending on the pathogen and host factors. Antimicrobial resistance is an emerging concern that complicates treatment and increases morbidity.
\section{Causes}
This condition happens because of fecal-oral transmission (genotypes 1, 2 through contaminated water) or consumption of undercooked pork or game meat (genotypes 3, 4). HEV replicates in hepatocytes, causing hepatocellular injury and tissue death. Following exposure, the pathogen invades host tissues and replicates, triggering an immune response that contributes to the symptoms. The severity of infection depends on pathogen virulence, host immune status, and timely treatment. Infection begins with pathogen entry through a susceptible portal (respiratory tract, gastrointestinal tract, skin, mucous membranes). The pathogen must overcome host defenses including physical barriers, innate immunity, and adaptive immune responses. Microbial virulence factors such as toxins, adhesins, and capsules enable the pathogen to establish infection and evade immune clearance. Host factors including age, nutritional status, immune competence, and comorbidities influence susceptibility and disease severity.
\section{Risk Factors}
Close contact with infected individuals Compromised immune system (HIV, immunosuppressive therapy, malnutrition) Extremes of age (very young or elderly) Travel to endemic regions Poor sanitation and hygiene practices Lack of vaccination Exposure to contaminated food or water Healthcare exposure (nosocomial infections) Occupational exposure (healthcare workers, laboratory personnel) Crowded living conditions

\begin{itemize}[noitemsep]
  \item Close contact with infected individuals
  \item Compromised immune system (HIV, immunosuppressive therapy, malnutrition)
  \item Extremes of age (very young or elderly)
  \item Travel to endemic regions
  \item Poor sanitation and hygiene practices
  \item Lack of vaccination
  \item Exposure to contaminated food or water
  \item Healthcare exposure (nosocomial infections)
  \item Occupational exposure (healthcare workers, laboratory personnel)
  \item Crowded living conditions
\end{itemize}
\section{Signs and Symptoms}
Symptoms of acute hepatitis E include jaundice, malaise, anorexia, abdominal pain, hepatomegaly, fever, and elevated transaminases, and are usually self-limited over 4--6 weeks. A distinctive feature is severe disease in pregnancy (genotypes 1, 2), with mortality of 20--30\% in the third trimester, high rates of fulminant liver failure, and adverse fetal outcomes (preterm delivery, stillbirth). Chronic HEV infection occurs in immunocompromised individuals (organ transplant recipients, HIV, hematologic malignancy) with genotypes 3, 4, leading to rapidly progressive cirrhosis. Fever and chills Fatigue and malaise Localized pain and inflammation at infection site Swelling, redness, and warmth (local infection) Cough and respiratory symptoms (respiratory infections) Nausea, vomiting, and diarrhea (gastrointestinal infections) Headache and body aches Skin rash or lesions Enlarged lymph nodes Systemic symptoms in severe cases (hypotension, organ dysfunction)

\begin{itemize}[noitemsep]
  \item Fever and chills
  \item Fatigue and malaise
  \item Localized pain and inflammation at infection site
  \item Swelling, redness, and warmth (local infection)
  \item Cough and respiratory symptoms (respiratory infections)
  \item Nausea, vomiting, and diarrhea (gastrointestinal infections)
  \item Headache and body aches
  \item Skin rash or lesions
  \item Enlarged lymph nodes
  \item Systemic symptoms in severe cases (hypotension, organ dysfunction)
\end{itemize}
\section{Progression and Stages}
The incubation period is the time between pathogen exposure and onset of symptoms, during which the pathogen replicates within the host. The prodromal phase involves early, nonspecific symptoms such as fever, fatigue, and malaise before characteristic symptoms appear. During the acute phase, hallmark signs and symptoms of the specific infection develop fully. The convalescent phase involves gradual resolution of symptoms as the immune system clears the infection. Some infections may progress to chronic or latent stages if the pathogen is not fully eliminated.
\section{Complications}
\begin{itemize}[noitemsep]
  \item Sepsis and septic shock
  \item Organ-specific complications (pneumonia, meningitis, endocarditis)
  \item Abscess formation requiring drainage
  \item Chronic infection (HIV, hepatitis B/C, tuberculosis)
  \item Reactive arthritis or post-infectious syndromes
  \item Antimicrobial resistance complicating treatment
  \item Disseminated infection in immunocompromised hosts
  \item Multi-organ failure in severe cases
\end{itemize}
\section{Diagnosis}
Doctors diagnose this based on serum anti-HEV IgM and HEV RNA detection (PCR). Laboratory diagnosis includes pathogen identification through culture, PCR testing, or antigen detection. Serological tests can confirm exposure or immune response to the pathogen. Detailed history including exposure, travel, and vaccination status Physical examination focusing on the affected system Complete blood count with differential Microbiological culture of blood, sputum, urine, or wound PCR and nucleic acid amplification tests for rapid pathogen detection Antigen detection tests Serological testing for antibodies (IgM, IgG) Imaging studies (chest X-ray, CT, ultrasound) Gram stain and microscopy of clinical specimens Antimicrobial susceptibility testing to guide therapy

\begin{itemize}[noitemsep]
  \item Detailed history including exposure, travel, and vaccination status
  \item Physical examination focusing on the affected system
  \item Complete blood count with differential
  \item Microbiological culture of blood, sputum, urine, or wound
  \item PCR and nucleic acid amplification tests for rapid pathogen detection
  \item Antigen detection tests
  \item Serological testing for antibodies (IgM, IgG)
  \item Imaging studies (chest X-ray, CT, ultrasound)
  \item Gram stain and microscopy of clinical specimens
  \item Antimicrobial susceptibility testing to guide therapy
\end{itemize}
\section{Prevention}
A recombinant HEV vaccine (Hecolin) is licensed in China. Hand hygiene with soap and water or alcohol-based sanitizer Vaccination according to recommended schedules Safe food handling and water purification Use of personal protective equipment in healthcare settings Safe sex practices including condom use Mosquito avoidance (nets, repellents, insecticides) Respiratory etiquette (covering coughs and sneezes) Isolation of infected individuals when indicated Prophylactic antibiotics or antivirals for high-risk exposures Travel precautions and pre-travel consultations

\begin{itemize}[noitemsep]
  \item Hand hygiene with soap and water or alcohol-based sanitizer
  \item Vaccination according to recommended schedules
  \item Safe food handling and water purification
  \item Use of personal protective equipment in healthcare settings
  \item Safe sex practices including condom use
  \item Mosquito avoidance (nets, repellents, insecticides)
  \item Respiratory etiquette (covering coughs and sneezes)
  \item Isolation of infected individuals when indicated
  \item Prophylactic antibiotics or antivirals for high-risk exposures
  \item Travel precautions and pre-travel consultations
\end{itemize}
\section{Treatment Options}
Treatment typically involves supportive for acute hepatitis E

\begin{itemize}[noitemsep]
  \item Ribavirin for chronic HEV (12 weeks)
  \item Reduction of immunosuppression is first-line for transplant recipients.
  \item Antibiotics for bacterial infections (targeted or empiric)
  \item Antiviral medications for viral infections
  \item Antifungal therapy for fungal infections
  \item Antiparasitic medications for parasitic infections
  \item Supportive care: fluids, rest, nutrition, fever control
  \item Hospitalization for severe cases requiring intravenous therapy
  \item Oxygen therapy and respiratory support if needed
  \item Surgical drainage of abscesses when necessary
  \item Pain management and symptom relief
  \item Treatment of complications and underlying conditions
\end{itemize}
\section{Living with It}
\begin{itemize}[noitemsep]
  \item Complete the full course of prescribed medications
  \item Get adequate rest and maintain hydration during recovery
  \item Monitor for worsening symptoms that require medical attention
  \item Practice good hygiene to prevent spread to others
  \item Follow up with your healthcare provider as recommended
  \item Gradually resume normal activities as symptoms improve
\end{itemize}
\section{Outlook}
\begin{itemize}[noitemsep]
  \item Most people recover well for acute HEV in immunocompetent hosts
  \item Chronic HEV is treatable with ribavirin.
\end{itemize}
